\section{Simulation Results and Benchmarking}
\label{sec:results}

\begin{table*}[t]
  \centering
  \caption{Core disturbance-family comparison for the benchmark study.
    RMSE values are reported in mHz and computed as family means of
    scenario means. The uncertainty term is the mean within-scenario Monte
    Carlo standard deviation for the winning estimator. $\Delta$ is the RMSE gap
    between the winner and the runner-up.}
  \label{tab:family_summary}
  \scriptsize
  \setlength{\tabcolsep}{2.8pt}
  \renewcommand{\arraystretch}{1.1}
  \begin{tabular}{l c l c l c l}
    \toprule
    Family & Count & RMSE winner & RMSE [mHz] & Runner-up & $\Delta$ [mHz] & Lowest $T_{\mathrm{trip}}$ [s] \\
    \midrule
    IEEE ramps & 9 & SOGI-PLL & $19.4 \pm 6.4$ & Koopman & 2.7 & 0 (multiple) \\
    Mag.\ steps & 7 & EKF & $0.388 \pm 0.143$ & RA-EKF & 1.45 & 0 (multiple) \\
    Modulation & 3 & RA-EKF & $17.8 \pm 21.4$ & Koopman & 13.7 & 0 (multiple) \\
    Phase jumps & 3 & SOGI-FLL & $174 \pm 67.5$ & PLL & 4.03 & ZCD: 0.0194 \\
    OOB interf. & 1 & EKF & $13.6 \pm 11.1$ & RA-EKF & 7.35 & EKF, RA-EKF: 0 \\
    IBR harmonics & 3 & EKF & $70.5 \pm 27.3$ & TFT & 4.13 & TFT, Prony: 0.00947 \\
    IBR ringdown & 5 & UKF & $313 \pm 23.2$ & EKF & 6.67 & ESPRIT: 0.0989 \\
    IBR multi-event & 1 & PLL & $173.5 \pm 66.7$ & UKF & 2.74 & IpDFT: 0.0439 \\
    \bottomrule
  \end{tabular}
\end{table*}

\subsection{No Single Estimator Dominates the Standard Scenario Families}

The standard disturbance families do not yield a single ranking
(Table~\ref{tab:family_summary}). SOGI-PLL gives the lowest mean RMSE on IEEE
ramp cases (19.4\,mHz). EKF dominates the magnitude-step family
(0.388\,mHz) and out-of-band interference (13.6\,mHz), and RA-EKF leads the
modulation family (17.8\,mHz). Zero-trip-risk ties are common in ramps,
magnitude steps, and
modulation, which means those families separate the estimators mainly by
accuracy rather than by protection-oriented risk. The margins are often small:
the ramp-family gap between SOGI-PLL and Koopman is only 2.7\,mHz, while the
gap between EKF and RA-EKF in magnitude steps is 1.45\,mHz.

The scenario-level winners show the same fragmentation. Across the 34
benchmark scenarios, nine different estimators win at least one scenario on
RMSE, and no estimator wins more than 10 of them. Even inside the ramp
family, SOGI-PLL wins six of the nine RMSE scenarios while RA-EKF takes the
other three. The magnitude-step family is the clearest exception: EKF wins all
seven magnitude-step scenarios and the family mean as well.

Table~\ref{tab:family_summary} also distinguishes a winner from a decisive
winner. In the magnitude-step family, EKF's 1.45\,mHz lead over RA-EKF is about
one order of magnitude larger than the winner's mean within-scenario spread
(0.143\,mHz), so the aggregate advantage is meaningful. By contrast, the ramp,
modulation, phase-jump, harmonic, ringdown, and out-of-band families all have
family-level margins that are smaller than or comparable to the same spread.
Those families should therefore be read as closely contested rather than
exhibiting stable single-method dominance.

Trip-risk is even less discriminating in the clean standard cases. In 23 of the
34 scenarios, the minimum trip-risk is zero, often
shared by 10 to 13 estimators in the ramp and magnitude-step families. In those
cases, trip-risk is useful mainly as a pass/fail screen, while RMSE and peak
error carry most of the ranking information.

Phase jumps behave differently. SOGI-FLL gives the lowest mean RMSE across the
phase-jump family (174\,mHz), but ZCD gives the lowest mean trip-risk
($0.0194$\,s). The protection and accuracy metrics therefore point to
different designs. In the $60^\circ$ phase-jump case, ZCD keeps trip-risk below
the leading low-RMSE methods, but only with much larger error excursions. Even
within the standard test set, RMSE and trip-risk are not interchangeable
objectives.

\subsection{Composite IBR Stress Reorders the Ranking}

The broadest change appears when the disturbance becomes composite. The hardest
benchmark scenarios, measured by median RMSE across estimators, are the IBR
multi-event sequence (826\,mHz), the IEEE and NERC $60^\circ$ phase jumps
(675\,mHz and 637\,mHz), and the noisy IBR ringdown scenarios
(474--528\,mHz). These are precisely the cases where the usual standard test
intuition starts to fail.

The scale change is itself important. The easiest family in
Table~\ref{tab:family_summary}, magnitude steps, stays below 1\,mHz, whereas
the IBR multi-event scenario reaches 826\,mHz median RMSE across estimators and
the noisy ringdown family sits near half a hertz. The benchmark therefore spans
roughly three orders of magnitude in error severity: the median RMSE of the
IBR multi-event case is about $1.55\times10^3$ times that of the single-tone
case. That spread explains why rankings inferred from clean standard tests do
not transfer directly to composite IBR stress.

The IBR multi-event sequence is the clearest example. PLL gives the lowest RMSE
of 173.5\,mHz. UKF trails by only 2.74\,mHz on the family mean. EKF remains
nearly tied in the scenario readout at 176.4\,mHz while reducing trip-risk to
$0.0874$\,s. IpDFT yields the lowest trip-risk ($0.0439$\,s) but at the cost
of a much larger peak error (2550\,mHz). The
ranking therefore depends on the deployment objective: minimizing average error,
minimizing time above the trip threshold, or bounding transient severity.

The same pattern appears in the other IBR-specific families. EKF gives the
lowest mean RMSE across the harmonic family (70.5\,mHz), but TFT and Prony
give the lowest mean trip-risk there ($0.00947$\,s). In the ringdown and noise
family, UKF gives the lowest mean RMSE (313\,mHz), whereas ESPRIT gives the
lowest mean trip-risk ($0.0989$\,s) at a far higher computational cost. These
results argue against a single global ranking. Scenario family is part of
the problem definition, not a \emph{post hoc} qualifier.

The scenario-level results make that split sharper. Among the 12 scenarios with a unique
trip-risk winner, 9 have a different winner than the RMSE leader. Those
disagreements concentrate in the $60^\circ$ phase jumps and in the harmonic,
ringdown, and multi-event IBR cases. The harmonic family is especially
instructive: TFT wins two of the three individual harmonic scenarios, and
RA-EKF wins the third. Yet EKF still achieves the best family-average RMSE
because it stays consistently close to the scenario leaders across all three
stress levels.

Composite stress also changes which metric carries the practical decision. In
magnitude steps and clean ramps, many leading estimators already tie at zero
trip-risk, so RMSE is the useful separator. In the multi-event, phase-jump, and
ringdown families, RMSE, trip-risk, and peak excursion split across different
estimators. That metric split, not just the change in rank order, is the core
reason a single scalar ranking fails to describe protection-relevant
behavior.

\subsection{Global Rankings Are Useful Only as Secondary Evidence}

Family-by-family reading is the main analysis, but the global summaries still
help when they are interpreted carefully. The interpretation depends on the
aggregation rule. Mean-rank summaries across the 34 benchmark scenarios favor
TFT on RMSE ($4.38$) and peak error
($4.32$), with RA-EKF close behind on both ($4.62$ and $4.59$). Simple
averaging over the 34 benchmark scenarios instead favors RA-EKF and EKF,
with mean RMSE values of 116\,mHz and 123\,mHz, followed by UKF at 131\,mHz.
That difference is informative rather than contradictory: mean rank rewards
repeated near-front finishes, whereas raw averaging is driven more strongly by
the hardest disturbance families.

Consistency statistics make that distinction concrete. EKF is the most frequent
accuracy leader, with 10 outright RMSE wins, 19 RMSE top-three finishes, and
20 peak-error top-three finishes across the 34 benchmark scenarios. TFT wins
only two RMSE scenarios, yet it reaches the RMSE top three in 16 scenarios and
the peak-error top three in 17, which explains its strong mean-rank position.
Trip-risk favors a different group again: IpDFT and ESPRIT enter the trip-risk
top three in 25 and 24 scenarios, respectively. Global summaries are therefore
best read as screening tools, not as substitutes for family-level evidence.

The MAE versus CPU Pareto front contains EKF, RA-EKF, SOGI-FLL, SOGI-PLL, and
ZCD, indicating that neither accuracy nor cost alone determines the admissible
set. As secondary evidence, the study
also includes ANOVA and Kruskal-Wallis family tests on RMSE, MAE, peak error,
trip-risk, settling time, and CPU time; both tests reject the null hypothesis
of equal family-level distributions for all six metrics (all $p<10^{-16}$).
These tests are coarse
and do not replace the scenario-family analysis, but they support the claim
that the observed ranking differences are not confined to one metric.

ZCD remains on the MAE versus CPU Pareto front because its cost is negligible,
not because its transient accuracy remains competitive under phase jumps or
multi-event stress. The practically relevant low-cost region is therefore
narrower than the raw Pareto set suggests.

For deployment decisions, the recommendation remains scenario-dependent.
SOGI-PLL is the strongest
choice when ramp RMSE is the main objective, EKF is strongest on magnitude-step
and out-of-band families, RA-EKF is strongest on modulation, and PLL is the
lowest-RMSE option in the composite IBR event. None of those statements can be
promoted to a universal recommendation, because the lowest trip-risk in the
same family may belong to a different method. That conditional reading is
stronger than a flat ranked list because it keeps both the disturbance family
and the operating objective in view.

\subsection{Cost Keeps Fast Embedded Methods in the Discussion}

CPU cost changes the ranking again. Figure~\ref{fig:dashboard_mega} gathers
representative time-domain cases together with the global summaries. EKF and
RA-EKF occupy the best global accuracy region with mean CPU costs of
$0.813\,\mu$s and $1.61\,\mu$s per sample. SOGI-PLL and SOGI-FLL remain much
cheaper ($0.0867\,\mu$s and $0.0348\,\mu$s) and therefore stay relevant for
embedded protection even when they lose some accuracy. Window-based methods
such as TFT and IpDFT improve global trip-risk relative to some loop methods,
but they do so at about $3$--$4\,\mu$s per sample. ESPRIT reaches the
trip-risk top three almost as often as IpDFT, but at $2838\,\mu$s per sample
instead of $3.68\,\mu$s. Koopman and ESPRIT therefore occupy a substantially
higher-cost region, with mean costs of $390\,\mu$s and $2838\,\mu$s, without a matching
global RMSE gain. The time-domain panels illustrate the same conclusion from a different
way: the estimator that looks best in the ramp case is not the one that best
manages the multi-event or ringdown case.

Taken together, the results suggest two practically relevant operating
regions. EKF and RA-EKF define the low-error middle-cost region: they are not
the cheapest estimators, but they remain competitive across the widest range of
families and give the strongest aggregate balance once composite IBR stress is
included. SOGI-PLL and SOGI-FLL define the low-cost embedded region: they keep
sub-microsecond cost while remaining credible on ramps and selected phase-jump
conditions. Window-based methods such as TFT and IpDFT remain useful when lower
trip-risk justifies extra compute, but the more expensive reduced-order methods
do not show a broad enough error advantage in this benchmark to justify routine
deployment.

\begin{figure*}[!t]
  \centering
  \includegraphics[width=\textwidth,height=0.835\textheight,keepaspectratio]{Fig2_Mega_Dashboard.png}
  \caption{Representative benchmark views from the study.
    Panels (a)--(d) and (f) show single-scenario responses, each summarized
    over the same 60 Monte Carlo runs used throughout the paper; panels (e),
    (g), and (h) summarize all 34 benchmark scenarios. (a) Phase-jump
    response. (b) Ramp response. (c) IBR multi-event response. (d) Ringdown
    response. (e) RMSE distributions by estimator family. (f) Harmonic-stress
    response. (g) Global mean MAE versus CPU time. (h) Global mean trip-risk
    versus CPU time. Together the panels show why the benchmark should be
    interpreted by disturbance family and operating objective rather than by a
    single scalar score.}
  \label{fig:dashboard_mega}
\end{figure*}
